\input{../tex-inputs/latex-leseansicht-vorspann}

\section[Rainer Maria Rilke an Arthur Schnitzler, {[}Anfang?{]} April 1896]{L04534 Rainer Maria Rilke an Arthur Schnitzler, {[}Anfang?{]} April 1896}
\nopagebreak\mylabel{L04534v}
\rehead{ }\normalsize\beginnumbering\briefempfaengerindex{Schnitzler, Arthur@\textsc{Schnitzler, Arthur}!zzzRilke, Rainer Maria@\emph{von Rainer Maria Rilke}!1896-04-102@{{[}Anfang?{]} April 1896}|(be}
\toendnotes[C]{\smallbreak\pagebreak[2]}
\correspDesc{Versand  durch Rainer Maria Rilke im Zeitraum [Anfang?] April 1896 in Prag
\newline{}Erhalt  durch Arthur Schnitzler in Wien}\toendnotes[C]{\smallbreak}
\Standort{CUL, Schnitzler, B 84.}
\physDesc{Brief, 1 Blatt, 2 Seiten, 1716 Zeichen
\newline{}Handschrift: schwarze Tinte, deutsche Kurrent
\newline{}Schnitzler: mit Bleistift Vermerk: »\textsc{Rilke}« }
\buchAbdrucke{\weitereDrucke{\emph{Rainer Maria Rilke und Arthur Schnitzler. Ihr Briefwechsel.}. Mit Anmerkungen herausgegeben von Heinrich Schnitzler In: \emph{Wort und Wahrheit}, Jg. 13, Nr. 4, April 1958, S. 282–298, hier S. 283.} }\toendnotes[C]{\smallbreak}
\pstart
           \centering{}{\pb}Prag\oindex{Prag@\textbf{Prag}|pw}, im \label{K_L04534-55v}\edtext{April 1896}{\lemma{\textnormal{\emph{April 1896}}}\Cendnote{\textnormal{Mutmaßlich zeitgleich mit dem Widmungsexemplar vom XXXX Auszeichnungsfehler: Dokument L04539 nicht gefunden versandt.}}}\label{K_L04534-55}.\pend
           
\pstart
           \centering{}\textcolor{gray}{\textbf{\textit{Prag II. Wassergasse 15\textsuperscript{B} I.\oindex{Vodičkova 15B@\textbf{Vodičkova 15B}|pw}}}}\pend
           
\pstart{}Hochwerter Meiſter,\pend\vspace{0.5em}
\pstart
           Das erſte Heft\pwindex{Rilke, Rainer Maria 4.\,12.\,1875 Prag – 29.\,12.\,1926 Montreux@\textsc{Rilke, Rainer Maria} (4.\,12.\,1875 Prag – 29.\,12.\,1926 Montreux)!Lieder dem Volke geschenkt@\strich\emph{Lieder dem Volke geschenkt}|pwv} meines \textsc{Volks-Gratis-Unternehmens} »\textsc{\uline{Wegwarten\pwindex{Wegwarten@\emph{Wegwarten}|pw}}}«, das ich mit Liedern im Volkston erfüllt habe, hat viel Beifall
               gefunden. – Ich habe mir{ }ſeinerzeit auch erlaubt, Ihnen ein{ }ſolches Liederheft
               zu \label{K_L04534-1v}\edtext{überſenden}{\lemma{\textnormal{\emph{übersenden}}}\Cendnote{\textnormal{Das erste Heft\pwindex{Rilke, Rainer Maria 4.\,12.\,1875 Prag – 29.\,12.\,1926 Montreux@\textsc{Rilke, Rainer Maria} (4.\,12.\,1875 Prag – 29.\,12.\,1926 Montreux)!Lieder dem Volke geschenkt@\strich\emph{Lieder dem Volke geschenkt}|pwkv} war am 2. 1. 1896 erschienen. Ein
                  Schnitzler gewidmetes Exemplar ist nicht überliefert.}}}\label{K_L04534-1}. Dieſes neue (II.) »\textsc{Wegwarten}«-Heft\pwindex{Wegwarten@\emph{Wegwarten}|pw} dürfte zweifelsfrei Ihr Intereſſe
               in höherem Maße beanſpruchen, da es eine kleine »\textsc{\uline{moderne Scene\pwindex{Rilke, Rainer Maria 4.\,12.\,1875 Prag – 29.\,12.\,1926 Montreux@\textsc{Rilke, Rainer Maria} (4.\,12.\,1875 Prag – 29.\,12.\,1926 Montreux)!Jetzt und in der Stunde unseres Absterbens. Scene@\strich\emph{Jetzt und in der Stunde unseres Absterbens. Scene}|pwv}}}« bringt.\pend
           
\pstart
           Es liegt mir viel daran, das Urthe{[}i{]}l des von mir hochgeſchätzten
               Dichters der »\textsc{\uline{Liebelei\pwindex{Schnitzler, Arthur 15.\,5.\,1862 Wien – 21.\,10.\,1931 ebd.@\textsc{Schnitzler, Arthur} (15.\,5.\,1862 Wien – 21.\,10.\,1931 ebd.), \emph{Schriftsteller, Mediziner}!Liebelei. Schauspiel in drei Akten@\strich\emph{Liebelei. Schauspiel in drei Akten}|pw}}}« zu vernehmen. –\pend
           
\pstart
           Darf ich noch ein paar Worte über das Weſen der »\textsc{Wegwarten\pwindex{Wegwarten@\emph{Wegwarten}|pw}}« anfügen. Leicht erinnern Sie{ }ſich, verehrteſter Herr Doctor, des Vorwortes\pwindex{Rilke, Rainer Maria 4.\,12.\,1875 Prag – 29.\,12.\,1926 Montreux@\textsc{Rilke, Rainer Maria} (4.\,12.\,1875 Prag – 29.\,12.\,1926 Montreux)!Wort nur@\strich\emph{Ein Wort nur}|pwv}, welches das 1. Heft\pwindex{Rilke, Rainer Maria 4.\,12.\,1875 Prag – 29.\,12.\,1926 Montreux@\textsc{Rilke, Rainer Maria} (4.\,12.\,1875 Prag – 29.\,12.\,1926 Montreux)!Lieder dem Volke geschenkt@\strich\emph{Lieder dem Volke geschenkt}|pwv} einleitete. – Dem Volke{ }ſind auch die billigſten
                     Ausgaben zu theuer, hieß es dort: »\label{K_L04534-2v}\edtext{\textsc{Wenn es zwei Kreuzer{ }ſind
                           und die Frage heißt: Buch oder Brot? Brot werden{ }ſie wählen. Und wollt ihrs
                           verargen? Drum wollt ihr Allen geben, so \uline{gebt}.\pwindex{Rilke, Rainer Maria 4.\,12.\,1875 Prag – 29.\,12.\,1926 Montreux@\textsc{Rilke, Rainer Maria} (4.\,12.\,1875 Prag – 29.\,12.\,1926 Montreux)!Wort nur@\strich\emph{Ein Wort nur}|pwv}}}{\lemma{\textnormal{\emph{Wenn … gebt.}}}\Cendnote{\textnormal{Die betreffende Stelle lautet: »Und wenn es zwei
                                    Kreuzer sind, und die Frage
                                    heisst: Buch oder Brot? Brot werden sie wählen; wollt ihr’s verargen? Wollt ihr also Allen geben, — so \so{gebt}! —«.}}}\label{K_L04534-2}\pwindex{Rilke, Rainer Maria 4.\,12.\,1875 Prag – 29.\,12.\,1926 Montreux@\textsc{Rilke, Rainer Maria} (4.\,12.\,1875 Prag – 29.\,12.\,1926 Montreux)!Wort nur@\strich\emph{Ein Wort nur}|pwv}« (S. [4])\pend
           
\pstart
           Das war der Grundgedanke, der mich bei der »Wegwarten\pwindex{Wegwarten@\emph{Wegwarten}|pw}«-Gründung leitete. {\pb}Von
                  \textsc{nächstem
                     Hefte\pwindex{Deutsch-moderne Dichtungen@\emph{Deutsch-moderne Dichtungen}|pwv}} an werden die »\textsc{Wegwarten\pwindex{Wegwarten@\emph{Wegwarten}|pw}}« Organ eines Bundes »\textsc{Moderner}\textsc{Fantasie Künstler}«, oder wie ich den Bund nennen werde. Die
               gemeinſchaftliche Idee der Mitglieder iſt: Modernes Schaffen, Unterwerfung unter die
               Macht der »Stimmung«, der \uline{intimen}\textsc{fantasie}vollen Sti{\geminationm}ung! –\pend
           
\pstart
           Mitglieder{ }ſind \textsc{Dichter} und auch bildende \textsc{Künstler}. Wenn Sie, Herr \textsc{Doctor},
               dieſer Idee, dieſem Leitmotiv nicht ferne{ }ſtehen, dann werde ich Sie von dem
               Fortſchreiten meines Planes unterrichten. Jedes Mitglied erhielte dann 60–100 \textsc{Wegwarten}nummern\pwindex{Wegwarten@\emph{Wegwarten}|pw} (frei) zu freier Vertheilung{\dotssix}\pend
           
\pstart
           Doch ich ermüde Sie mit dieſer Langathmigkeit, ohne mich vorher vergewiſſert zu
               haben, ob mein Plan Ihr Intereſſe hat.\pend
           
\pstart
           Ich wünſchte es von ganzer Seele!\pend
           
\pstart
           In größter Verehrung{\\[\baselineskip]}\spacefill\mbox{RenéMRilke}\pend
           \leftskip=0em{}\selectlanguage{ngerman}\endnumbering\briefempfaengerindex{Schnitzler, Arthur@\textsc{Schnitzler, Arthur}!zzzRilke, Rainer Maria@\emph{von Rainer Maria Rilke}!1896-04-012@{{[}Anfang?{]} April 1896}|)be}\mylabel{L04534h}
\begin{anhang}
\end{anhang}\newcommand{\dateiname}{L04534}\newcommand{\titel}{Rainer Maria Rilke an Arthur Schnitzler, [Anfang?] April 1896}\newcommand{\editorInnen}{Herausgegeben von Jahnke, SelmaMüller, Martin Anton}\input{../tex-inputs/latex-leseansicht-abspann}
